\scanpage{15}\mspage{14}
En résumé, en termes de $\unsure{\mathcal{F}|\mathscr{F}}$, c'est de la donnée de deux \uline{relations} \uline{d'ordre}
\[
\boxed{\underset{\overbrace{\scriptstyle\text{rel. d'ordre dans }\mathcal{F}}}{\leq,\quad\ll}\quad\text{telles que}\quad F\leq G \Rightarrow F\ll G}
\tag*{(1.30)}
\]
que je vais partir.

En termes de celles-ci, on \unsure{definit|définit} \deleted{$\unsure{M|\mathcal{F}}$}
\[
M\subset\mathcal{F}
\]
par (1.20), avec les relations $\underset{M}{\leq},\ \underset{M}{\ll}$ induites par $\underset{\mathcal{F}}{\leq},\ \underset{\unsure{\mathcal{F}|\overline{f}}}{\ll}$, puis
\[
\mathcal{L}\subset M
\]
par (1.2\unsure{6|5}). L'application
\[
\mathcal{F}\hookrightarrow\unsure{\mathcal{P}|\mathrm{P}}(M)\qquad
\deleted{\unsure{X}}F\longmapsto\widetilde F=\unsure{\mathrm{Dep}|\mathrm{Def}}(F)
\tag*{($*$)}
\]
est \inserted{alors} \unsure{définie|definie} par
\[
\widetilde F=\{X\in M\mid X\leq F\}
\tag*{(1.31)}
\]
et
\[
\mathcal{F}\hookrightarrow\mathrm{Fig}(M)\qquad
\deleted{\unsure{X}}F\longmapsto\mathrm{Fig}_{M}(F)
\tag*{($**$)}
\]
par (1.4) :
\[
\mathrm{Fig}_{M}(F)=\{\widetilde X\mid X\in\widetilde F\}\subset\mathcal{P}(M)
\tag*{\deleted{(1.32)}}
\]
avec
\[
\widetilde X=\{Y\in M\mid Y\leq X\}\quad.
\]

Ainsi, $\mathcal{F}$ est interprété comme un ens. de parties dans $M$ via ($*$), et $\leq$ \unsure{s'interprète|s'interprête} simplement par l'inclusion
\[
F\leq G\Longleftrightarrow\widetilde F\subset\widetilde G
\tag*{(1.32)}
\]
